\documentclass[twocolumn]{article}    % two-column because cheat sheets fit like that
\usepackage[utf8]{inputenc}
\usepackage{amsmath}    % allows math to happen
\usepackage{marvosym}
\usepackage{textcomp}
\usepackage{geometry}   % allows you to change the margins and headers/footers, and more
\usepackage{array}
\usepackage{wrapfig}
\usepackage{multirow} 
\usepackage{graphicx}   % allows for inclusion of graphics or several types, i think
\usepackage{float}      % SUPER  IMPORTANT. lets you tell LaTeX where to put your graphics, as opposed to where it feels
\usepackage{wrapfig}    % Lets text wrap around figures, more relevant for reports really
\usepackage{caption}
\usepackage{amssymb}    
\usepackage{ulem}       % allows strikethrough and some
\usepackage[table]{xcolor}  % enhances tables by allowing colors to happen. notably, as background/fill colors
\usepackage{bondgraph}  % allows bond graphs
\usepackage{caption}    % lots of control of captions and caption geometries
\usepackage{comment} % allows block comments to occur
\usepackage{blindtext}

\geometry{letterpaper, total={4.8in, 2.8in}}      

\begin{document}
\pagestyle{empty}
\footnotesize                      % small because, it's a cheat sheet after all, why not cram it

\begin{center}
    Levi Barker \hspace{35pt} U0959362
\end{center}
\vspace{-10pt}
\underline{General Knowledge}
\begin{itemize}
    \item \(det(A)\neq0\), invertible, full rank, lin. ind. cols/rows, nullity=0, non-singular
    \item Zero-state is "synonym" for TF thinking
    \item D, U, and L matrices show eigs on diag 
    \item A J Matrix has 1 block for each grade-1 eigenvector. Ex: 3x3 w/ multiplicity 3 results in one J block as the entire matrix
    \item Eigvector order in \(V\) and \(J\) match 
    \item nullity\((A-\lambda_iI)\) = \# of J-blocks for \(\lambda_i\)
    \item Basis is lin. ind. cols of \(A\), null space is remaining cols of \(A\)
    \item Z transform wants func. of \(z^{-1}\) not \(z\) 
    \item A square matrix satisfies it's own \(\Delta(\lambda)\)
    \item SS models are equiv. if same eigenvalues. Zero-state equiv. if same TF (matrix)
\end{itemize}
\begin{equation*}
    \hspace{-10pt} G(s) = C(Is-A)^{-1}B+D \hspace{5pt} AQ = Q\Lambda \hspace{5pt} AV=VJ
\end{equation*}
\begin{equation*}
    \overbrace{y_3(t) = \alpha y_1(t)}^\text{Homogeneity} \hspace{5pt}\text{if} \hspace{5pt}
    \begin{cases}
        x_3(0) = \alpha x_1(0)\\
        u_3(t) = \alpha u_1(t)
    \end{cases}
\end{equation*}
\begin{equation*}
    \hspace{-7.5pt}
    \overbrace{y_3(t) = y_1(t) + y_2(t)}^\text{Additivity} \text{if}
    \begin{cases}
        x_3(0) = x_1(0) + x_2(0) \\
        u_3(t) = u_1(t) + u_2(t)
    \end{cases}
\end{equation*}
\begin{equation*}
    \hat{A}_{\text{comp.}} = \mathcal{C}^{-1} A \mathcal{C} \hspace{15pt} \text{follow through in sys.}
\end{equation*}
\begin{equation*}
    \underbrace{\begin{bmatrix}
        -\alpha_1 & -\alpha_2 & -\alpha_3 & -\alpha_4 \\
        1 & 0 & 0 & 0 \\
        0 & 1 & 0 & 0 \\
        0 & 0 & 1 & 0
    \end{bmatrix}}_{ \hspace{5pt}
    \Delta(\lambda) = 
    \lambda^4 + \alpha_1 \lambda^3 + \alpha_2 \lambda^2 + \alpha \lambda + \alpha_4}
\end{equation*}
\begin{equation*}
    \hspace{-10pt} (A-I\lambda)v_i = q_i \hspace{2.5pt} \Longleftrightarrow \hspace{2.5pt} v_i \text{ is gen. eigvec of } A
\end{equation*}
\begin{align*}
    y[k+2] + 5y[k+1] + 10y[k] &=  7u[k] \\
    z^2Y[z] +5zY[z]+10Y[z]&=7U[z]
\end{align*}
\underline{General Solutions}\\
\begin{equation*}
    x(t) = e^{At}x(0) + \int_0^te^{A(t-\tau)} Bu(t)d\tau
\end{equation*}
\begin{equation*}
    x[k] = A^kx[0] \hspace{10pt} \text{  for y, plug x into y eq.}
\end{equation*}
\begin{equation*}
    x(t) = e^{Q\Lambda Q^{-1}t}x(0) = Qe^{\Lambda t} Q^{-1} x(0)
\end{equation*}
\underline{Functions of Matrices}\\
\begin{equation*}
    \begin{matrix}
        e^{At}= \\
        \\
        \text{(applies to} \\
        \text{J blocks)}
    \end{matrix}
    \begin{bmatrix}
        e^{\lambda t} & t e^{\lambda t} & \frac{t^2}{2!} e^{\lambda t} & \frac{t^3}{3!} e^{\lambda t}\\
        0 & e^{\lambda t} & t e^{\lambda t} & \frac{t^2}{2!} e^{\lambda t} \\
        0 & 0 & e^{\lambda t} & t e^{\lambda t} \\
        0 & 0 & 0 & e^{\lambda t}
    \end{bmatrix} \hspace{5pt} 
\end{equation*}
\begin{equation*}
    e^A = 1 + A^2/2! + A^3/3! + A^4/4! +...
\end{equation*}
\begin{equation*}
    sin(A) = A - A^3/3! +A^5/5! - A^7/7!+ ...
\end{equation*}
\begin{equation*}
    cos(A) = 1 - A^2/2! + A^4/4! - A^6/6! +... 
\end{equation*}
\begin{equation*}
    f(\hat{A}) = Q^{-1}f(A)Q \hspace{15pt} f(A) = Q^{-1}f(\hat{A})Q
\end{equation*}
\begin{equation*}
    \hspace{-5pt} \text{if } A = 
    \begin{bmatrix}
        A_1 & 0\\
        0& A_2
    \end{bmatrix}
    \rightarrow f(A) = 
    \begin{bmatrix}
        f(A_1) & 0\\
        0 & f(A_2)\\
    \end{bmatrix}
\end{equation*}
\begin{equation*}
    \begin{matrix}
        J^k = \\
        \text{(single}\\
        \text{block)}
    \end{matrix} 
    \begin{bmatrix}
        \lambda^k & k\lambda^{k-1} & \frac{k(k-1)}{2!} \lambda^{k-2}\\
        0 & \lambda^k & k\lambda^{k-1} \\
        0 & 0 & \lambda^k \\
    \end{bmatrix}
\end{equation*}
\noindent
\underline{Polynomials}\\
\(\Delta(\lambda)\) has all eigvalues. \(\psi (\lambda)\) only includes the single largest J block for each eigvalue  Ex:
\begin{equation*}
    \hspace{-15pt}
    \begin{bmatrix}
        \lambda_1 & 1 & 0 &0\\
        0 &  \lambda_1 & 1 & 0\\
        0 & 0 & \lambda_1 & 0\\
        0 & 0 & 0 & \lambda_2
    \end{bmatrix}  
    \begin{cases}
        \Delta(\lambda) = (\lambda-\lambda_1)^3(\lambda-\lambda_2)\\
        \psi (\lambda) = (\lambda-\lambda_1)^3(\lambda-\lambda_2)
    \end{cases}
\end{equation*}
\begin{equation*}
    \begin{bmatrix}
        \lambda_1 & 1 & 0 &0\\
        0 &  \lambda_1 & 0 & 0\\
        0 & 0 & \lambda_1 & 0\\
        0 & 0 & 0 & \lambda_1
    \end{bmatrix} 
    \begin{cases}
        \Delta (\lambda) = (\lambda-\lambda_1)^4\\
        \psi (\lambda) = (\lambda-\lambda_1)^2
    \end{cases}
\end{equation*}

\noindent
\underline{Realizations (method two)} \hspace{15pt} (for i inputs)\\
\begin{equation*}
    G(s) = \frac{\begin{bmatrix}
        \beta_{11}s^3 & \beta_{12}s^2 & \beta_{13}s & \beta_{14}\\
        \beta_{21}s^3 & \beta_{22}s^2 & \beta_{23}s & \beta_{24}
    \end{bmatrix}} {s^4 + \alpha_1s^3 + \alpha_2s^2 + \alpha_3 s +\alpha_4} + \overbrace{\begin{bmatrix}
        d_1\\ d_2
    \end{bmatrix}}^{G(\infty)}
\end{equation*}
\begin{equation*}
    A = \begin{bmatrix}
        -\alpha_1 & -\alpha_2 & -\alpha_3 & -\alpha_4 \\
        1 & 0  & 0 & 0 \\
        0 & 1 & 0 & 0 \\
        0 & 0 & 1 & 0\\
    \end{bmatrix}
\end{equation*}
\begin{equation*}
    B_i = \begin{matrix}
        [1 & 0]^T 
    \end{matrix} \hspace{30pt} D_i = \begin{matrix}
        [d_1 & d_2]^T
    \end{matrix}
\end{equation*}
\begin{equation*}
    C_i = \begin{bmatrix}
        \beta_{11} & \beta_{12} & \beta_{13} & \beta_{14}\\
        \beta_{21} & \beta_{22} & \beta_{23} & \beta_{24}
    \end{bmatrix}
\end{equation*}

Then 
\begin{equation*}
\begin{bmatrix}
    \Dot{x}_1 \\ \Dot{x}_2
\end{bmatrix} = \begin{bmatrix}
    A_1 & 0 \\ 0 & A_2
\end{bmatrix} \begin{bmatrix}
    x_1\\x_2
\end{bmatrix} + \begin{bmatrix}
    b_1 & 0 \\ 0 & b_2
\end{bmatrix} u
\end{equation*}
\begin{equation*}
    y = \begin{bmatrix}
        C_1 & C_2
    \end{bmatrix} \begin{bmatrix}
    x_1\\x_2
\end{bmatrix} + \begin{bmatrix}
    d_1 & d_2 
\end{bmatrix} u
\end{equation*}

\noindent
\underline{Stability}\\
BIBO: zero-state response \(\to G(0)*u(t)\). Stable iff every pole if OLHP/open unit circle\\
AS: all \(Re(\lambda_i)<0\) for C, or \(|\lambda_i|<1\) for D. And \(x_i(\infty) \to 0 \hspace{5pt} \forall i\)\\
MS: Allow \(Re(\lambda_i) = 0\) for C, or \(|\lambda_i|=1\) for D, only for simple roots (J-blocks sized 1)
    
\noindent
\underline{Controllable If}
\begin{itemize}
    \item \(\rho(\mathcal{C}) = \rho(\begin{matrix}
        [B & AB & ... & A^{n-1}B]
    \end{matrix}) = n\) 
    \item \(\rho(\begin{matrix}
        [A-\lambda I & B]
    \end{matrix}) = n \hspace{5pt} \forall \lambda_i\)
    \item Rows(\(B\)) aligned with bottom of each J-block are lin ind if repeated, else nonzero
\end{itemize}


\noindent
\underline{Observable If}
\begin{itemize}
    \item \(\rho(\mathcal{O}) = \rho\begin{matrix}
        [C & CA & ... & CA^{n-1}]^T 
    \end{matrix}=n\)
    \item \(\rho(\begin{matrix}
        [A-\lambda I & C]^T
    \end{matrix}) = n \hspace{5pt} \forall \lambda_i\) 
    \item Cols(\(C\)) aligned with left of each J-block are lin ind if repeated, else nonzero
\end{itemize}

\noindent
\underline{Kalman Decomposition}\\
\textbf{if} \(\rho(\mathcal{C}) = n_1<n\):\\
form \(P^{-1} = \begin{matrix}
        [q_1 & ... & q_{n_1} & ... q_n]
        \end{matrix}\) with  \(n_1\) lin ind cols of \(\mathcal{C}\), and rest to make invertible, then \(P^{-1} \hat{x} = x\) and 
        \begin{equation*}
            \begin{bmatrix}
                \hat{\Dot{x}}_C \\ \hat{\Dot{x}}_{\overline{C} }
            \end{bmatrix} = \begin{bmatrix}
                \hat{A}_C & \hat{A}_{12} \\ 
                0 & \hat{A}_{\overline{C}}
            \end{bmatrix} \begin{bmatrix}
                \hat{x}_C \\ \hat{x}_{\overline{C}}
            \end{bmatrix} + \begin{bmatrix}
                \hat{B}_C \\ 0
            \end{bmatrix} u
        \end{equation*}
        \begin{equation*}
            y = \begin{bmatrix}
                \hat{C}_C & \hat{C}_{\overline{C}}
            \end{bmatrix} \begin{bmatrix}
                \hat{x}_C \\ \hat{x}_{\overline{C}}
            \end{bmatrix} + Du
        \end{equation*}
        
        \noindent
        \textbf{if} \(\rho(\mathcal{O}) =n_2<n\):\\
        form \(P = \begin{matrix}
            [P_1 & ... & P_{n_2} & ... P_n]^T
        \end{matrix}\) with \(n_2\) lin ind rows of \(\mathcal{O}\), and rest to make invertible, then \(P^{-1} \hat{x} = x\) and 
        \begin{equation*}
            \begin{bmatrix}
                \hat{\Dot{x}}_O \\ \hat{\Dot{x}}_{\overline{O} }
            \end{bmatrix} = \begin{bmatrix}
                \hat{A}_O & 0 \\ 
                \hat{A}_{21} & \hat{A}_{\overline{O}}
            \end{bmatrix} \begin{bmatrix}
                \hat{x}_O \\ \hat{x}_{\overline{O}}
            \end{bmatrix} + \begin{bmatrix}
                \hat{B}_O \\ \hat{B}_{\overline{O}}
            \end{bmatrix} u
        \end{equation*}
        \begin{equation*}
            y = \begin{bmatrix}
                \hat{C}_O & 0
            \end{bmatrix} \begin{bmatrix}
                \hat{x}_O \\ \hat{x}_{\overline{O}}
            \end{bmatrix} + Du
        \end{equation*}
    
   
    
\noindent
\underline{Controllability After Sampling}\\
If \((A,B)\) is controllable, a sufficient\\ 
condition for \((A_d, B_d)\) controllability is:
\begin{equation*}
    \text{if } Re(\lambda_i) - Re(\lambda_j) = 0 \hspace{10pt} \lambda_i , \lambda_j \in A, \text{ then}
\end{equation*}
\begin{equation*}
    \text{pick } T: |Im(\lambda_i)| - |Im(\lambda_j)| \neq \frac{2\pi m}{T} \hspace{5pt} m \in \mathcal{N}_1
\end{equation*}



\noindent
\underline{CL State Feedback}
\begin{equation*}
    u = r - k \otimes x \hspace{10pt}  \Dot{x} = (A-Bk) x + Br
\end{equation*}
\begin{equation*}
    (A-Bk, B) \text{ is controllable iff } (A,B) \text{ is}
\end{equation*}
\begin{equation*}
    k_i = \hat{\alpha}_i - \alpha_i = \text{des }- \text{current}
\end{equation*}
\begin{equation*}
    x_{ss} = -(A-Bk)^{-1}Br \hspace{20pt} y_{ss} = Cx_{ss}
\end{equation*}
\begin{enumerate}
    \item Controllable? 
    \item Find \(A_{CL}= (A-B\otimes k)\) in terms of \(k_i\)
    \item \(\Delta_{CL}(s) = det(\lambda I-A_{CL}) \rightarrow \alpha_i\)
    \item Determine \(\alpha_i\) from \(\Delta_{des}(s)\)
    \item \(g_{dc}= C(A-Bk)^{-1}B+D \hspace{20pt} r = y/ g_{dc}\)
\end{enumerate}
\underline{Stabilization}
\begin{equation*}
    \hspace{-25pt} \begin{bmatrix}
        \hat{\Dot{x}}_C \\ \hat{\Dot{x}}_{\overline{C}} 
    \end{bmatrix} = \begin{bmatrix}
        \hat{A}_C - \hat{B}_C\hat{k}_C & \hat{A}_{12} -\hat{B}_C\hat{k}_{\overline{C}} \\
        0 & \hat{A}_{\overline{C}}
    \end{bmatrix} \begin{bmatrix}
        \hat{x}_C \\ \hat{x}_{\overline{C}} 
    \end{bmatrix} + \begin{bmatrix}
        \hat{B}_C \\ 0
    \end{bmatrix} r
\end{equation*}
\begin{equation*}
    \text{if } u = r - \begin{matrix}
        [k_C & k_{\overline{C}}]
    \end{matrix} \otimes \begin{matrix}
        [x_C & x_{\overline{C}}]
    \end{matrix}^T
\end{equation*}
Stabilizable: if uncontrollable states are already stable

\vspace{30pt}
\begin{center}
    Levi Barker \hspace{35pt} U0959362
\end{center}



%% NOT INCLUDED IN 3x5s
\newpage 
\begin{equation*}
    A^{-1} = \frac{A_{co}^T}{det(A)} \hspace{15pt} ||x||_k \triangleq \left( \Sigma_{i=1}^n |x_i|^k \right)^{1/k}
\end{equation*}


\noindent
\textbf{Method ONE}: \hspace{60pt} \(D = G(\infty)\)
\begin{equation*}
    G_{sp}(s) = \frac{1}{D(s)}
    \begin{matrix}
        [N_1s^{n-1} & N_2s^{n-2} & ... & N_n]\\
    \end{matrix} 
\end{equation*}
\begin{equation*}
    A = \begin{bmatrix}
        -\alpha_1I & -\alpha_2I & ... & -\alpha_{n-1}I & -\alpha_nI \\
        I & 0 & ... & 0 & 0 \\
        0 & I & ...& 0 & 0 \\
          &   & \ddots &  & \\
        0 & 0 & ... & I & 0\\
    \end{bmatrix}
\end{equation*}
\begin{equation*}
    B = \begin{matrix}
        [I & 0 & 0 & ... & 0]^T
    \end{matrix} \hspace{9pt} C =[N_1 + ... + N_n]
\end{equation*}
Will work, but not in fewest states.

\noindent
\textbf{Method TWO}: \hspace{60pt} \(D = G(\infty)\)
Generate SS models column by column
\begin{equation*}
    \begin{bmatrix}
        \Dot{x}_1 \\ \Dot{x}_2
    \end{bmatrix} = \begin{bmatrix}
        A_1 & 0\\ 0 & A_2
    \end{bmatrix}
    \begin{bmatrix}
        x_1 \\ x_2
    \end{bmatrix} + \begin{bmatrix}
        b_1 & 0\\ 0 & b_2
    \end{bmatrix} \begin{bmatrix}
        u_1 \\ u_2
    \end{bmatrix}
\end{equation*}
\begin{equation*}
    y = \begin{bmatrix}
        C_1 & C_2
    \end{bmatrix} \Vec{x} + G(\infty)\Vec{u}
\end{equation*}
Will work, in fewest possible states.


    \noindent
    \underline{LQR}
    \begin{equation*}
        \min(J) = \int_0^\infty (x^T(\tau)Qx(\tau) + u^T(\tau)Ru(\tau))d\tau
    \end{equation*}
    \(Q\) must be +ve semi-definite. \(x^T(t)Qx(t) \geq 0 \forall x(t) \neq 0\). Diags must be positive OR zero. If any off-diag. causes any \(\lambda=0\), not semi-definite.\\
    \(R\) must be positive definite. \(u^T(t)Ru(t) > 0 \forall u(t) \neq 0\). Diags must be greater than zero.
    \begin{equation*}
        u = -Kx \hspace{10pt} K = R^{-1} B^T P \text{ where } P \text{ solves:}
    \end{equation*}
    \begin{equation*}
        ARE: \hspace{10pt} A^TP + PA - PBR^{-1}B^TP + Q =0
    \end{equation*}
    LQR is iterative. Pick \(Q\) and \(R\), run sim with hard ic's and tweak to achieve desired behavior of system.
    
\newpage
\noindent
\underline{Full Kalman Decomposition}
    \begin{align*}
        \hspace{-5pt}
        \begin{bmatrix}
            \hat{\Dot{x}}_{CO}\\ 
            \hat{\Dot{x}}_{C\overline{O}}\\
            \hat{\Dot{x}}_{\overline{C}O}\\
            \hat{\Dot{x}}_{\overline{CO}}\\
        \end{bmatrix} =& \begin{bmatrix}
            \hat{A}_{CO} & 0 & \hat{A}_{13} & 0 \\
            \hat{A}_{21} & \hat{A}_{C\overline{O}} & \hat{A}_{23} & \hat{A}_{24} \\
            0 & 0 & \hat{A}_{\overline{C}O} & 0\\
            0 & 0 &\hat{A}_{43} & \hat{A}_{\overline{CO}}
        \end{bmatrix} \begin{bmatrix}
            \hat{x}_{CO}\\ 
            \hat{x}_{C\overline{O}}\\
            \hat{x}_{\overline{C}O}\\
            \hat{x}_{\overline{CO}}
        \end{bmatrix}\\
        &+ \begin{bmatrix}
            \hat{B}_{CO} & 
            \hat{B}_{C\overline{O}} &  0 & 0
        \end{bmatrix}^T u
    \end{align*}
    \begin{equation*}
        y = \begin{bmatrix}
            \hat{C}_{CO} & 0 & \hat{C}_{\overline{C}O} & 0 
        \end{bmatrix} \begin{bmatrix}
            \hat{x}_{CO}\\ 
            \hat{x}_{C\overline{O}}\\
            \hat{x}_{\overline{C}O}\\
            \hat{x}_{\overline{CO}}
        \end{bmatrix} + Du
    \end{equation*}
    
    \noindent
    \underline{Kalman Filter}\\
    Incorporate system and sensor noise as a benefit    
    \begin{enumerate}
        \item The discrete system model 
        \begin{align*}
            x_{t+1} =& Ax_t + Bu_t + m \hspace{5pt} &m \sim \mathcal{N}(0,M)\\
            z_{t+1} =& Cx_t + n          & n \sim  \mathcal{N}(0,N)
        \end{align*}
        
        \item Incorporate the past system behavior
        \begin{equation*}
            x_{t-1}  \sim \mathcal{N}(\hat{x}_{t-1}, \varepsilon_{t-1})
        \end{equation*}
        
        \item Do the process update (in code)
        \begin{align*}
            \hat{x}_t' =& A\hat{x}_{t-1} + Bu_{t-1}\\
            \varepsilon_t' =& A\varepsilon_{t-1}A^T + m
        \end{align*}
        
        \item Do the measurement update (in code)
        \begin{align*}
            K_t =& \varepsilon_t' C^T (C\varepsilon_t' C^T+ n)^{-1}\\
            \hat{x}_t =& \overbrace{\hat{x}_t'}^{\text{w/o sensors}} + \overbrace{K_t \underbrace{(z_t-C\hat{x}_t')}_{\text{innovation}} }^{\text{w/ sensors}}\\
            \varepsilon_t =& \varepsilon_t' - K_t C \varepsilon_t' \hspace{7.5pt} \text{// Update uncert.}
        \end{align*}
    \end{enumerate}
    \vspace{20pt}
    \begin{center}
        Levi Barker \hspace{35pt} U0959362
    \end{center}
\end{document}